% Oversells — ilmiy maqola shabloni (XeLaTeX + BibTeX).
% Article template with abstract, table and bibliography (compile with XeLaTeX).
\documentclass[11pt,a4paper]{article}

\usepackage[margin=2.5cm]{geometry}
\usepackage{fontspec}
\setmainfont{Noto Serif}
\setsansfont{Noto Sans}
\usepackage{babel}
\babelprovide[import=uz-Latn, main]{uzbek}
\setlength{\emergencystretch}{3em}
% babel'ning oʻzbek lokalida yoʻq nomlar / names missing from babel's Uzbek locale
\setlocalecaption{uzbek}{abstract}{Annotatsiya}
\setlocalecaption{uzbek}{table}{Jadval}
\setlocalecaption{uzbek}{ref}{Adabiyotlar}

\usepackage{amsmath, amssymb, amsthm}
\usepackage{booktabs}
\usepackage[hidelinks]{hyperref}

\newtheorem{teorema}{Teorema}

\title{Maktab matematikasida masalalar yechish usullari}
\author{Muallif Ismi\thanks{Toshkent shahri, 1-maktab. E-pochta: muallif@example.uz}}
\date{\today}

\begin{document}

\maketitle

\begin{abstract}
  Ushbu maqolada umumtaʼlim maktablarida kvadrat tenglamalarni oʻqitishning
  turli usullari koʻrib chiqiladi va ularning samaradorligi solishtiriladi.
  \par\medskip
  \noindent\textbf{Kalit soʻzlar:} kvadrat tenglama, diskriminant, Viyet teoremasi, metodika.
\end{abstract}

\section{Kirish}

Kvadrat tenglamalar algebra kursining asosiy mavzularidan biridir~\cite{alimov}.
Mavzuni oʻqitishda grafik, analitik va amaliy yondashuvlar qoʻllaniladi~\cite{polya}.

\section{Asosiy natija}

\begin{teorema}[Viyet]
  Agar $x_1$ va $x_2$ sonlari $x^2 + px + q = 0$ tenglamaning ildizlari boʻlsa, u holda
  \begin{equation}
    x_1 + x_2 = -p, \qquad x_1 x_2 = q.
  \end{equation}
\end{teorema}

\section{Tajriba natijalari}

\begin{table}[h]
  \centering
  \caption{Nazorat ishlari natijalari (oʻrtacha ball)}
  \begin{tabular}{lcc}
    \toprule
    Usul & Tajriba guruhi & Nazorat guruhi \\
    \midrule
    Analitik & 4,2 & 3,8 \\
    Grafik & 4,5 & 3,9 \\
    Amaliy masalalar & 4,6 & 4,0 \\
    \bottomrule
  \end{tabular}
\end{table}

\section{Xulosa}

Natijalar shuni koʻrsatadiki, amaliy masalalar orqali oʻqitish oʻquvchilarning
mavzuni oʻzlashtirishini yaxshilaydi.

\bibliographystyle{unsrt}
\bibliography{adabiyotlar}

\end{document}
